\documentclass[aspectratio=169, 11pt]{beamer}

% ═══════════════════════════════════════════════════════════════════
% CONFIGURACIÓN — estilo limpio, color sobrio (azul apagado)
% ═══════════════════════════════════════════════════════════════════
\usepackage[utf8]{inputenc}
\usepackage[T1]{fontenc}
\usepackage[spanish]{babel}
\usepackage{lmodern}
\usepackage{booktabs}
\usepackage{array}
\usepackage{xcolor}
\usepackage{colortbl}

\usetheme{default}
\useinnertheme{circles}

% Paleta: azul sobrio profesional + gris
\definecolor{azulin}{RGB}{45,82,130}
\definecolor{azulsuave}{RGB}{225,232,242}
\definecolor{gristexto}{RGB}{70,70,70}
\definecolor{verdeclaro}{RGB}{240,255,244}
\definecolor{rojoclaro}{RGB}{255,245,245}
\definecolor{azulclaro}{RGB}{235,248,255}

\setbeamercolor{frametitle}{fg=azulin}
\setbeamercolor{title}{fg=azulin}
\setbeamercolor{structure}{fg=azulin}
\setbeamercolor{itemize item}{fg=azulin}
\setbeamercolor{itemize subitem}{fg=gristexto}
\setbeamercolor{enumerate item}{fg=azulin}
\setbeamercolor{block title}{bg=azulin, fg=white}
\setbeamercolor{block body}{bg=azulsuave, fg=black}
\setbeamercolor{normal text}{fg=black}

\setbeamertemplate{navigation symbols}{}

\setbeamertemplate{footline}{%
    \hfill\insertframenumber\,/\,\inserttotalframenumber\hspace{0.4cm}\vspace{0.2cm}
}

\setbeamertemplate{frametitle}{%
    \vspace{0.3cm}
    {\color{azulin}\large\bfseries\insertframetitle}\par
    {\color{azulin}\rule{\linewidth}{1pt}}
}

\setbeamerfont{title}{size=\Large, series=\bfseries}

% ═══════════════════════════════════════════════════════════════════
\begin{document}

% ─── PORTADA ───
\begin{frame}[plain]
    \centering
    \vspace{0.5cm}
    {\footnotesize UNIVERSIDAD NACIONAL DE SAN ANTONIO ABAD DEL CUSCO}\\
    {\footnotesize Escuela Profesional de Ingeniería Informática y de Sistemas}\\
    {\footnotesize Curso: IF614 -- Ingeniería de Software I}\\[0.4cm]

    {\color{azulin}\rule{\linewidth}{1.2pt}}\\[0.4cm]
    {\Large\bfseries\color{azulin} SISTEMA INTELIGENTE DE RECOLECCIÓN}\\[0.15cm]
    {\Large\bfseries\color{azulin} DE RESIDUOS SÓLIDOS SEGREGADOS -- CUSCO}\\[0.3cm]
    {\large\bfseries\color{azulin} Entrega N.° 2 -- Documentación y sistema funcional}\\[0.4cm]
    {\color{azulin}\rule{\linewidth}{1.2pt}}\\[0.5cm]

    {\normalsize Metodología: SCRUM \quad$|$\quad Atributo del graduado: AG-C01}\\[0.3cm]
    {\normalsize Semestre 2026-I}\\[0.5cm]

    {\footnotesize
    \textbf{Integrantes:}\\
    Luicho Quispe, Jhoel Alex (224871) $|$
    Luna Ccapa, Carlos Willian (210178)\\
    Puma Condori, Richard Braulio (161809) $|$
    Huaracallo Arenas, Lino Zeynt (204798)\\[0.2cm]
    Cusco -- Perú, 2026}
\end{frame}

% ─── AGENDA ───
\begin{frame}{Contenido}
    \begin{enumerate}
        \item \textbf{Problema y objetivos} del proyecto
        \item \textbf{Alcance} y público objetivo
        \item \textbf{Metodología SCRUM}: roles y Product Backlog
        \item \textbf{Sprints} realizados
        \item \textbf{Arquitectura y tecnologías}
        \item \textbf{Sistema construido}: backend, app móvil y dashboard web
        \item \textbf{Despliegue} en la nube y APK
        \item \textbf{Diagramas} (UML / Entidad-Relación)
        \item \textbf{Resultados} y \textbf{conclusiones}
    \end{enumerate}
\end{frame}

% ─── PROBLEMA Y OBJETIVOS ───
\begin{frame}{1 — Problema y Objetivos}
    \vspace{-0.3cm}
    \begin{columns}[T]
    \column{0.48\textwidth}
    \begin{block}{Problema}
        \footnotesize
        En el Cusco la recolección de residuos presenta:
        \begin{itemize}
            \item Rutas no optimizadas, sin trazabilidad
            \item Escasa comunicación con la población
            \item Baja cultura de segregación
            \item Ausencia de reportes para decidir
        \end{itemize}
    \end{block}

    \column{0.48\textwidth}
    \begin{block}{Objetivo general}
        \footnotesize
        Diseñar e implementar un sistema inteligente de gestión de residuos sólidos segregados que optimice la recolección, mejore la comunicación con la ciudadanía y genere información para la toma de decisiones, aplicando SCRUM.
    \end{block}
    \end{columns}
    \vspace{0.15cm}
    \begin{block}{Objetivos específicos}
        \footnotesize
        Gestión de usuarios y zonas $\cdot$ catálogo y guía de segregación $\cdot$ monitoreo de rutas con ubicación $\cdot$ app móvil ciudadana $\cdot$ sistema de alertas $\cdot$ reportes y analítica.
    \end{block}
\end{frame}

% ─── ALCANCE ───
\begin{frame}{2 — Alcance y Público Objetivo}
    \vspace{-0.3cm}
    \begin{columns}[T]
    \column{0.5\textwidth}
    \begin{block}{Alcance}
        \footnotesize
        \begin{itemize}
            \item Administración de usuarios y zonas
            \item Catálogo de residuos y segregación
            \item Monitoreo de rutas
            \item App móvil ciudadana
            \item Alertas y reportes
        \end{itemize}
        \textit{\scriptsize Fuera de alcance: hardware IoT, pasarelas de pago, sistemas tributarios.}
    \end{block}

    \column{0.46\textwidth}
    \begin{block}{Público objetivo}
        \footnotesize
        \begin{itemize}
            \item \textbf{Ciudadanos:} rutas, alertas, reportar incidencias, segregar
            \item \textbf{Operadores:} ejecutar rutas, enviar ubicación
            \item \textbf{Administradores:} zonas, rutas, usuarios, reportes
        \end{itemize}
    \end{block}
    \end{columns}
\end{frame}

% ─── SCRUM: ROLES + BACKLOG ───
\begin{frame}{3 — Metodología SCRUM: Roles}
    \vspace{-0.2cm}
    \begin{block}{Equipo Scrum}
        \footnotesize
        \begin{table}
        \scriptsize
        \begin{tabular}{|l|l|l|}
        \hline
        \rowcolor{azulin}
        \textcolor{white}{\textbf{Rol}} & \textcolor{white}{\textbf{Integrante}} & \textcolor{white}{\textbf{Función}} \\
        \hline
        Scrum Master & Huaracallo Arenas, Lino Zeynt & Facilita el proceso, remueve impedimentos \\
        \hline
        Product Owner & Luicho Quispe, Jhoel Alex & Prioriza el backlog, criterios de aceptación \\
        \hline
        Developer (Full Stack) & Luna Ccapa, Carlos Willian & Backend, app móvil y dashboard \\
        \hline
        Developer (Full Stack) & Puma Condori, Richard Braulio & Backend, app móvil y dashboard \\
        \hline
        Developer (Full Stack) & Luicho Quispe, Jhoel Alex & Integración y despliegue \\
        \hline
        \end{tabular}
        \end{table}
        \vspace{0.1cm}
        \scriptsize Todos los integrantes participan como desarrolladores full stack; Lino asume además el rol de Scrum Master.
    \end{block}
    \begin{block}{Herramientas de colaboración}
        \footnotesize
        GitHub (repositorio y ramas) $\cdot$ Render (despliegue) $\cdot$ MongoDB Atlas (BD) $\cdot$ Expo (app móvil).
    \end{block}
\end{frame}

% ─── PRODUCT BACKLOG ───
\begin{frame}{3 — Product Backlog (Requisitos)}
    \vspace{-0.35cm}
    \begin{columns}[T]
    \column{0.56\textwidth}
    \begin{block}{Requisitos Funcionales}
        \scriptsize
        \begin{tabular}{|l|l|c|}
        \hline
        \rowcolor{azulin}
        \textcolor{white}{\textbf{ID}} & \textcolor{white}{\textbf{Requisito}} & \textcolor{white}{\textbf{Est.}}\\
        \hline
        RF-01 & Registro de ciudadanos & \checkmark \\
        RF-02 & Autenticación (JWT) & \checkmark \\
        RF-03/04 & Zonas + asignación (GPS) & \checkmark \\
        RF-05/06 & Tipos y clasificación de residuos & \checkmark \\
        RF-07/08/09 & Rutas + ubicación + gestión & \checkmark \\
        RF-10 & Consulta de rutas (móvil) & \checkmark \\
        RF-11 & Reporte de incidencias (GPS) & \checkmark \\
        RF-14 & Reporte de residuos por zona & \checkmark \\
        RF-12/13/15/16 & Alertas y reportes avanzados & $\sim$ \\
        \hline
        \end{tabular}
    \end{block}

    \column{0.42\textwidth}
    \begin{block}{Requisitos No Funcionales}
        \scriptsize
        \begin{itemize}
            \item Seguridad: bcrypt + JWT
            \item Disponibilidad: nube 24/7
            \item Usabilidad: interfaz simple
            \item Rendimiento: $<$ 3 s
            \item Portabilidad: APK Android
            \item Privacidad: Ley 29733
        \end{itemize}
    \end{block}
    \end{columns}
    \vspace{0.1cm}
    \centering{\scriptsize \checkmark\ completado \quad $\sim$\ parcial}
\end{frame}

% ─── SPRINTS ───
\begin{frame}{4 — Sprints Realizados}
    \vspace{-0.2cm}
    \begin{block}{3 Sprints · 2 semanas cada uno}
        \footnotesize
        \begin{table}
        \scriptsize
        \begin{tabular}{|l|p{4.2cm}|p{4.8cm}|}
        \hline
        \rowcolor{azulin}
        \textcolor{white}{\textbf{Sprint}} & \textcolor{white}{\textbf{Objetivo}} & \textcolor{white}{\textbf{Entregables}}\\
        \hline
        \textbf{1} & Autenticación y gestión de usuarios/zonas & API auth (JWT), modelos, zonas, detección por GPS \\
        \hline
        \textbf{2} & Núcleo operativo + app móvil & Rutas + ubicación, residuos, incidencias, pantallas móviles \\
        \hline
        \textbf{3} & Alertas, reportes, web y despliegue & Dashboard web, reportes, despliegue en la nube, APK \\
        \hline
        \end{tabular}
        \end{table}
    \end{block}
    \begin{block}{Ramas por historia (GitHub)}
        \scriptsize
        \texttt{feature/HU01-registro-ciudadanos} \dots \texttt{feature/HU08-dashboard-admin}\\
        Flujo: rama por HU $\rightarrow$ commit $\rightarrow$ Pull Request $\rightarrow$ merge a \texttt{main} $\rightarrow$ despliegue.
    \end{block}
\end{frame}

% ─── ARQUITECTURA Y TECNOLOGÍAS ───
\begin{frame}{5 — Arquitectura y Tecnologías}
    \vspace{-0.3cm}
    \begin{columns}[T]
    \column{0.5\textwidth}
    \begin{block}{Arquitectura (cliente–servidor)}
        \scriptsize
        \begin{itemize}
            \item \textbf{Clientes:} app móvil (Android) + dashboard web
            \item \textbf{Servidor:} API REST (Node/Express) en Render
            \item \textbf{Datos:} MongoDB Atlas (nube)
            \item Comunicación HTTPS/JSON, autenticación JWT
        \end{itemize}
    \end{block}

    \column{0.48\textwidth}
    \begin{block}{Tecnologías}
        \scriptsize
        \begin{itemize}
            \item Backend: \textbf{Node.js + Express}
            \item Base de datos: \textbf{MongoDB} (Mongoose)
            \item App: \textbf{Expo / React Native} (TypeScript)
            \item Web: HTML + JS (servido por Express)
            \item Mapas: OpenStreetMap / Leaflet
            \item Integración: API RENIEC (DNI)
        \end{itemize}
    \end{block}
    \end{columns}
    \vspace{0.1cm}
    \begin{block}{Flujo}
        \centering\scriptsize
        App móvil / Web $\;\longrightarrow\;$ API REST (Render) $\;\longrightarrow\;$ MongoDB Atlas
    \end{block}
\end{frame}

% ─── SISTEMA: BACKEND ───
\begin{frame}{6 — Sistema: Backend (API REST)}
    \vspace{-0.3cm}
    \begin{columns}[T]
    \column{0.48\textwidth}
    \begin{block}{Modelos de datos}
        \footnotesize
        \begin{itemize}
            \item Usuario (roles, zona, ubicación)
            \item Zona (polígono geográfico)
            \item Ruta (estado, ubicación en vivo)
            \item Incidente, Residuo, Alerta
        \end{itemize}
    \end{block}
    \begin{block}{Seguridad}
        \footnotesize
        Contraseñas con \textbf{bcrypt} $\cdot$ sesión con \textbf{JWT} $\cdot$ control de acceso por rol.
    \end{block}

    \column{0.48\textwidth}
    \begin{block}{Endpoints principales}
        \scriptsize
        \begin{itemize}
            \item \texttt{/auth/register}, \texttt{/auth/login}
            \item \texttt{/zonas}, \texttt{/zonas/detect} (GPS)
            \item \texttt{/rutas}, \texttt{/rutas/:id/ubicacion}
            \item \texttt{/incidentes}, \texttt{/residuos}
            \item \texttt{/usuarios}, \texttt{/alertas}, \texttt{/dni/:dni}
        \end{itemize}
    \end{block}
    \begin{block}{Detección de zona}
        \scriptsize
        Consulta geoespacial \textbf{punto-en-polígono} para asignar la zona del ciudadano por su ubicación.
    \end{block}
    \end{columns}
\end{frame}

% ─── SISTEMA: APP MÓVIL + WEB ───
\begin{frame}{6 — Sistema: App Móvil y Dashboard Web}
    \vspace{-0.3cm}
    \begin{columns}[T]
    \column{0.5\textwidth}
    \begin{block}{App móvil (Android / APK)}
        \footnotesize
        \begin{itemize}
            \item Login y registro (mapa + GPS)
            \item Autocompletar datos por \textbf{DNI} (RENIEC)
            \item Inicio: rutas y estado
            \item \textbf{Mapa: camiones en vivo}
            \item Reporte de incidencias con foto/ubicación
            \item Guía de segregación $\cdot$ Perfil
        \end{itemize}
    \end{block}

    \column{0.46\textwidth}
    \begin{block}{Dashboard web (administración)}
        \footnotesize
        \begin{itemize}
            \item Resumen / indicadores
            \item Gestión de usuarios y zonas
            \item Gestión de rutas y estados
            \item Gestión de incidencias
            \item Catálogo de residuos
        \end{itemize}
        \scriptsize Servido por el mismo backend (misma URL).
    \end{block}
    \end{columns}
    \vspace{0.1cm}
    \centering{\scriptsize \textit{[Insertar capturas reales de la app y el dashboard]}}
\end{frame}

% ─── DESPLIEGUE ───
\begin{frame}{7 — Despliegue en la Nube y APK}
    \vspace{-0.3cm}
    \begin{columns}[T]
    \column{0.5\textwidth}
    \begin{block}{Despliegue}
        \footnotesize
        \begin{itemize}
            \item \textbf{Backend + Web:} Render (en línea 24/7)
            \item \textbf{Base de datos:} MongoDB Atlas
            \item \textbf{Repositorio:} GitHub (main + ramas)
            \item \textbf{App:} APK Android instalable
        \end{itemize}
    \end{block}

    \column{0.48\textwidth}
    \begin{block}{Funciona desde cualquier red}
        \footnotesize
        El sistema deja de depender de una PC local: la app y el dashboard se conectan a un backend en la nube con URL permanente.
    \end{block}
    \end{columns}
    \vspace{0.1cm}
    \begin{block}{Estado de la Entrega 2}
        \footnotesize
        Sistema \textbf{funcional} (supera el 50\% requerido): backend + app + dashboard desplegados. Documentación: Capítulo I, II y bosquejo del III.
    \end{block}
\end{frame}

% ─── DIAGRAMAS ───
\begin{frame}{8 — Diagramas (UML / Entidad-Relación)}
    \vspace{-0.2cm}
    \begin{block}{Diagramas elaborados}
        \footnotesize
        \begin{itemize}
            \item \textbf{Casos de uso} — actores: ciudadano, operador, administrador
            \item \textbf{Clases / dominio} — Usuario, Zona, Ruta, Incidente, Residuo, Alerta
            \item \textbf{Entidad-Relación} — colecciones y relaciones en MongoDB
            \item \textbf{Secuencia} — login y reporte de incidencia
            \item \textbf{Componentes} — clientes, API y base de datos
        \end{itemize}
    \end{block}
    \centering{\scriptsize \textit{[Insertar aquí los diagramas — disponibles en docs/DIAGRAMAS.md]}}
\end{frame}

% ─── RESULTADOS ───
\begin{frame}{9 — Resultados}
    \vspace{-0.3cm}
    \begin{columns}[T]
    \column{0.5\textwidth}
    \begin{block}{Logros}
        \footnotesize
        \begin{itemize}
            \item Sistema de 3 capas funcional y desplegado
            \item 11 de 16 requisitos funcionales completos
            \item App móvil instalable (APK)
            \item Integración con API RENIEC (DNI)
        \end{itemize}
    \end{block}

    \column{0.48\textwidth}
    \begin{block}{Impacto (AG-C01)}
        \footnotesize
        \begin{itemize}
            \item \textbf{Ambiental:} promueve la segregación
            \item \textbf{Social:} participación ciudadana
            \item \textbf{Legal:} Ley 29733, Ley 27314, NTP 900.058
        \end{itemize}
    \end{block}
    \end{columns}
    \vspace{0.1cm}
    \begin{block}{Pendiente}
        \footnotesize
        Notificaciones push de cercanía/retraso (RF-12/13) y reportes avanzados (RF-15/16) — previstos para la siguiente etapa.
    \end{block}
\end{frame}

% ─── CONCLUSIONES ───
\begin{frame}{Conclusiones}
    \vspace{-0.3cm}
    \begin{enumerate}
        \item \textbf{Sistema funcional desplegado}\\
        \footnotesize Backend (Node/Express + MongoDB), app móvil (Expo) y dashboard web, en la nube — supera el 50\% requerido para la Entrega 2.
        \vspace{0.1cm}
        \item \textbf{SCRUM aplicado}\\
        \footnotesize Product Backlog priorizado, 3 sprints, ramas por historia de usuario y trabajo colaborativo en GitHub.
        \vspace{0.1cm}
        \item \textbf{Cobertura de requisitos}\\
        \footnotesize 11/16 RF completos; los módulos esenciales (usuarios, zonas, rutas, residuos, incidencias) operativos.
        \vspace{0.1cm}
        \item \textbf{Responsabilidad ética y legal}\\
        \footnotesize Cifrado de datos, manejo responsable de información personal (Ley 29733) y alineación con normas ambientales.
        \vspace{0.1cm}
        \item \textbf{Próximos pasos}\\
        \footnotesize Notificaciones push, reportes avanzados y refinamiento del dashboard para la entrega final.
    \end{enumerate}
\end{frame}

% ─── CIERRE ───
\begin{frame}[plain]
    \centering
    \vspace{2.5cm}
    {\color{azulin}\rule{0.5\linewidth}{1pt}}\\[0.5cm]
    {\Large\bfseries\color{azulin} Gracias por su atención}\\[0.4cm]
    {\normalsize ¿Preguntas?}\\[0.4cm]
    {\color{azulin}\rule{0.5\linewidth}{1pt}}
\end{frame}

\end{document}
